% CLAUDE REPORT::WORSTCASE-0001::2026-11-05
% !TEX program = xelatex
\documentclass[11pt]{article}
\usepackage[margin=1in]{geometry}
\usepackage{amsmath,amssymb,amsthm,xurl}
\newtheorem{theorem}{Theorem}
\newtheorem{proposition}[theorem]{Proposition}
\newcommand{\SAT}{\mathrm{SAT}}
\newcommand{\UNSAT}{\mathrm{UNSAT}}
\newcommand{\status}[1]{\textbf{[#1]}}
\title{CLAUDE WORSTCASE-0001\\[2pt]\large Worst-case minimum derivation cost $\overline W_R(N)$:\\ what the conditional bridge $\SAT\in P\Rightarrow \overline W_R\le \mathrm{poly}$ can and cannot be}
\author{CLAUDE, scheduled run}
\date{2026-11-05}
\begin{document}\maketitle
\noindent\textbf{Status.} P versus NP remains OPEN. No lower bound or transfer theorem is claimed. Additive; no schedules, allocations or root \texttt{main.tex} changed.

\section{Setup}
A system $R$ is a sound derivation relation $F\vdash_R\bot$ with a cost (size; or time/space of a verified derivation). $\overline W_R(N)=\max_{F\in\UNSAT,|F|\le N}\min_{\pi}\mathrm{cost}_R(\pi)$. The minimum is taken over \emph{all} sound derivations of that one $F$, with the best derivation given free.

\section{Results}
\begin{proposition}[Soundness is where the cost lives] \status{PROVED/EXACT}
If a derivation of $\bot$ from $F$ is allowed to be any object whose correctness is not independently checked, the cost is $O(1)$ (``output UNSAT''). Hence a meaningful $R$ must have derivations checkable in time polynomial in $|F|+|\pi|$ (Cook--Reckhow). Then $\overline W_R(N)$ is exactly the optimal proof length of the proof system $R$, and ``free best sequence'' is the standard definition of proof size. Charging nothing for search is therefore not a relaxation; it is the definition.
\end{proposition}

\begin{proposition}[Dichotomy for the bridge] \status{PROVED/EXACT}
Let $B(R)$ denote ``$\SAT\in P\Rightarrow \overline W_R(N)\le\mathrm{poly}(N)$''.
\begin{enumerate}
\item If $R$ has an unconditional superpolynomial lower bound (resolution, tree/DAG; bounded-depth Frege; cutting planes; polynomial calculus), then $\overline W_R$ is superpolynomial, so $B(R)$ is logically equivalent to $\SAT\notin P$, i.e. to $P\ne NP$. The bridge is not weaker than the target.
\item If no superpolynomial lower bound for $R$ is known (Frege, Extended Frege), then $B(R)$ is the statement ``$P=NP\Rightarrow R$ is polynomially bounded''. This is a genuine intermediate statement, but it asserts nothing about hardness: it is a \emph{simulation} claim.
\item $\overline W_R$ polynomial for \emph{every} sound checkable $R$ is impossible unless $NP=coNP$ fails to be refuted; conversely, $\overline W_R$ superpolynomial for \emph{all} $R$ is equivalent to $NP\ne coNP$ (Cook--Reckhow), which is strictly stronger than $P\ne NP$ (it implies it).
\end{enumerate}
\end{proposition}
\emph{Proof.} (1) is the contrapositive reading: a true conclusion makes the implication equivalent to the negation of its hypothesis. (3) is Cook--Reckhow 1979. $\square$

\begin{proposition}[What the bridge needs in case 2] \status{CANDIDATE / OPEN}
Proving $B(\mathrm{Frege})$ requires turning a polynomial SAT algorithm into polynomial-size Frege refutations. The natural route is formalizing the algorithm's correctness inside Frege (as bounded arithmetic does for $\mathsf{VNC}^1$-style theories). This needs the correctness of the algorithm to be provable in a theory whose proofs translate to short Frege proofs ($\mathsf{PV}_1$ translates to EF). Pure existence of a polynomial algorithm does \emph{not} give provability of its correctness; this is the known gap between $P=NP$ and $NP=coNP$-with-short-proofs. A transfer theorem would have to bridge it and I know of none.
\end{proposition}

\section{Consequences for the five assignments}
\begin{itemize}
\item \textbf{Lower-bound lead:} the weakest sufficient bridge is not weaker for any $R$ with known bounds; for $R$ with unknown bounds the bridge is a simulation conjecture, and a superpolynomial $\overline W_R$ for Frege would already be a new result (best known size bounds are still near-quadratic).
\item \textbf{Hard families / red team:} any family claimed hard must survive Frege/EF simulation of the algorithm; hardness for resolution-like $R$ is classical (Haken; Ben-Sasson--Wigderson) and says nothing about $P$.
\item \textbf{Time--space auditor:} for time/space cost, Atserias--Dalmau resolution space and size--space tradeoffs are known; they are model-specific and do not transfer to RAM time.
\item \textbf{Classification rule (use the world ledger):} ``$\overline W_R$ polynomial for all $R$ iff $NP=coNP$''. A candidate lemma equal to this is a reformulation.
\end{itemize}

\section{Literature check (2026-11-05)}
de~Rezende--Engstr\"om--Ghannane--Risse, ECCC TR26-078 (2026): superpolynomial length bounds for \emph{tree-like semantic Frege with bounded line size}; the authors state it is unclear whether such tree-like systems relate to DAG-like Frege, so this is not a general Frege bound. A community page reports best known Frege/EF size bounds still quadratic (unverified against a primary source). Khaniki--Pich--Sokolov, CCC 2026, prover--adversary games for Frege: not read in detail. Status of transfer: directly transferable theorem: none; analogy: bounded-line-size restrictions as a hardness lever, needs a bridge to DAG-like Frege.

\section{Corpus}
Reviewed: \texttt{AGENTS.md}, world ledger and its README, ONETAPE-0001, Grok multisign notes (the multisign DAG is explicitly \textbf{not} PROVED/EXACT for SAT there). The directive file referenced in the task is not present in the repository; this report does not reconstruct it.

\noindent\textbf{Bottom line.} P versus NP remains OPEN. Next target: decide whether any natural intermediate system between DAG-like resolution and Frege admits a superpolynomial bound surviving the algorithm-formalization argument of Proposition 3.
\end{document}
% END_OF_VERIFIED_REPORT::CLAUDE::WORSTCASE-0001
